\section*{Chapter \ref{ch:b2:three-equation-model} --- The Three-Equation Model}

\begin{solution}{exo:b2:three-equation-model:1}
% figdata: rule_dec, rule_summary
The 1970s, 4.74 points below on average. The mean absolute distance was 2.90 points in 1960--79 and 1.40 in 1987--2007.
\end{solution}

\begin{solution}{exo:b2:three-equation-model:2}
% figdata: recent_summary
Inflation peaked at 9.1 per cent in June 2022, when the funds rate was 1.21 per cent. The rate first exceeded half a point in May 2022
and peaked at 5.33 per cent in August 2023.
\end{solution}

\begin{solution}{exo:b2:three-equation-model:3}
% code: exercises/solutions.py ws_ps
$1.2 - 4u = 0.8$ gives $u = 0.1$, 10 per cent.
\end{solution}

\begin{solution}{exo:b2:three-equation-model:4}
% code: exercises/solutions.py taylor1993
$4 + 0.5 \times (-1) + 0.5 \times (4 - 2) + 2 = 6.5$ per cent.
\end{solution}

\begin{solution}{exo:b2:three-equation-model:5}
% code: exercises/solutions.py decay, rule_coef
MR: $x = -(\pi - \pi^{\mathrm T})$. The gap left is $1/(1 + 1) = 0.5$. The rule is $r - r^* = 1/(1 \times (1 + 1)) \times (\pi - \pi^{\mathrm T}) = 0.5(\pi - \pi^{\mathrm T})$.
\end{solution}

\begin{solution}{exo:b2:three-equation-model:6}
% code: exercises/solutions.py first_step
Inflation this period is $2 + 0.5 \times 2 = 3$. The bank chooses $\pi = 2 + 1 \times \tfrac23 = 2.667$ and $x = -0.667$ for next period, and sets $r = 2 + 0.667/0.5 = 3.333$; the nominal rate is $3.333 + 3 = 6.333$.
A demand shock that has already happened has raised inflation, and the bank must now bring it down: it pushes output below
equilibrium next period.
\end{solution}

\begin{solution}{exo:b2:three-equation-model:7}
% code: exercises/solutions.py zlb
Desired real rate $1 + \tfrac43 \times (0.5 - 2) = -1$, nominal $-1 + 0.5 = -0.5$. The rate set is 0, the real rate $-0.5$. Next period's gap would be $-0.5 \times (-1 - 1) = 1.0$ without the bound and $-0.5 \times (-0.5 - 1) = 0.75$ with it.
\end{solution}

\begin{solution}{exo:b2:three-equation-model:8}
% code: exercises/solutions.py first_step, cumulated_loss
The fraction left is $1/(1 + 0.25/2) = 0.889$: first-period inflation 4.667 rather than 4, real rate 3.333 rather than 6, output gap $-0.667$ rather than $-2$. The
cumulated loss is still 6.00: the same total, spread over more periods.
\end{solution}

\begin{solution}{exo:b2:three-equation-model:9}
% code: exercises/solutions.py residuals_by_decade
Means: 1960s $-0.49$, 1970s $-4.74$, 1980s 0.88, 1990s $-0.21$, 2000s $-2.13$. Mean absolute values: 1.05, 4.74, 2.52, 1.20, 2.29. The 1970s stand out: the mean equals the mean absolute value, so the rate was below the rule
in every quarter of the decade.
\end{solution}

\begin{solution}{exo:b2:three-equation-model:10}
Minimising $(\pi_t - \pi^{\mathrm T})^2 + \lambda x_t^2$ subject to $\pi_t = \pi_{t-1} + \alpha x_t$: the first-order condition in $x_t$ is $2\alpha(\pi_t - \pi^{\mathrm T}) + 2\lambda x_t = 0$, the MR curve. Substituting $x_t = (\pi_t - \pi_{t-1})/\alpha$: $\pi_t - \pi^{\mathrm T} + (\lambda/\alpha^2)(\pi_t - \pi_{t-1}) = 0$, so
$(\pi_t - \pi^{\mathrm T})(1 + \lambda/\alpha^2) = (\lambda/\alpha^2)(\pi_{t-1} - \pi^{\mathrm T})$ and $\pi_t - \pi^{\mathrm T} = (\pi_{t-1} - \pi^{\mathrm T})/(1 + \alpha^2/\lambda)$. Then $x_t = -(\alpha/\lambda)(\pi_t - \pi^{\mathrm T}) = -\dfrac{\alpha}{\lambda + \alpha^2}(\pi_{t-1} - \pi^{\mathrm T})$, and the IS curve gives $r_{t-1} - r^* = -x_t/a = \dfrac{\alpha}{a(\lambda + \alpha^2)}(\pi_{t-1} - \pi^{\mathrm T})$.
\end{solution}

\begin{solution}{exo:b2:three-equation-model:11}
% code: exercises/solutions.py stability_factor
The derived rule has $\phi - 1 = \alpha/(a(\lambda + \alpha^2)) > 0$, and $1 - a\alpha(\phi - 1) = 1 - \alpha^2/(\lambda + \alpha^2) = \lambda/(\lambda + \alpha^2)$, between 0 and 1. With Brenland's parameters the upper bound is $1 + 2/(0.5 \times 0.5) = 9$, far above the derived 2.33.
\end{solution}

\begin{solution}{exo:b2:three-equation-model:12}
Along the path $x_t = -(\alpha/\lambda)(\pi_t - \pi^{\mathrm T})$ and the Phillips curve gives $\alpha x_t = \pi_t - \pi_{t-1}$. Summing over $t \geq 1$: $\alpha\sum x_t = \lim \pi_t - \pi_0 = \pi^{\mathrm T} - \pi_0 = -g_0$, so the cumulated loss $-\sum x_t = g_0/\alpha$, whatever $\lambda$.
\end{solution}

\begin{solution}{pb:b2:three-equation-model:1}
% code: exercises/solutions.py BRENLAND, ws_ps, first_step, periods_below, cumulated_loss; listings/three_eq.py
\textbf{Part I.}
1. Workers obtain a lower real wage when unemployment is higher, since losing a job is costlier and the outside option worse.
2. $1 - 0.20 = 0.8$.
3. $1.1 - 5u = 0.8$, $u = 6$ per cent.
4. Wage setters claim $1.1 - 0.2 = 0.9$, more than the 0.8 prices allow, so wages and prices chase each other and inflation rises.
5. $1.1 - 5u = 0.75$, $u = 7$ per cent.
\textbf{Part II.}
6. $\pi_1 = 5 + 0.5x_1$.
7. Minimising $(\pi - 2)^2 + 0.5x^2$ subject to the Phillips curve: $2 \times 0.5(\pi - 2) + 2 \times 0.5x = 0$, so $x = -(\pi - 2)$.
8. $5 + 0.5x = 2 - x$: $x = -2$, $\pi = 4$.
9. $-2 = -0.5(r - 2)$, so $r = 6$ per cent.
10. $6 + 5 = 11$ per cent.
11. $r - 2 = \tfrac{0.5}{0.5 \times (0.5 + 0.25)}(\pi - 2) = \tfrac43(\pi - 2)$.
\textbf{Part III.}
12. $1/(1 + 0.25/0.5) = 2/3$.
13. $\pi_2 = 2 + 3 \times \tfrac49 = 3.33$; $x_2 = -1.33$.
14. $2 + \tfrac43 \times 2 = 4.67$ per cent.
15. $3(2/3)^t < 0.5$ for $t \geq 5$: inflation is 2.40 in period 5.
16. Period 0: 5.00, 6.00, 11.00; period 1: 4.00, 4.67, 8.67, as computed.
17. The nominal rate is $r + \pi = 2 + \pi + \tfrac43(\pi - 2)$, so it rises by $1 + \tfrac43 = 2.33$ points per point of inflation: more than one, as the \omterm{def:b2:three-equation-model:taylor-principle}{Taylor principle} requires.
\textbf{Part IV.}
18. $3/0.5 = 6$ per cent of a period's output.
19. $6/3 = 2$.
20. The gap left is $1/(1 + 0.25/2) = 0.889$: $\pi_1 = 4.67$, $x_1 = -0.25 \times 2.67 = -0.67$.
21. Still 6.
22. Each point of disinflation needs $1/\alpha$ points of negative output gap through the Phillips curve; the weight decides only how fast the bank
takes them.
23. It doubles, to 4.
24. No: the rule asks for a real rate of $0.5 + 4 = 4.5$ and a nominal rate of 9.5 per cent, far above zero. The bound binds when inflation and $r^*$ are both low.
25. The real rate starts at 6 per cent (nominal 11) and falls back toward 2 as inflation closes a third of its gap each period; the
cumulated output loss is 6 per cent of a period's output and the sacrifice ratio $1/\alpha = 2$.
\end{solution}
